\section{D3 — Base teórica}

Uma fila no vertiporto pode ter como unidade uma aeronave que solicita uma posição de pouso ou um passageiro que chega à inspeção. A definição dessa unidade determina o significado da espera. Em uma fila estável, a relação de \citet{little1961}, $L=\lambda W$, associa o número médio de unidades presentes à taxa média de unidades admitidas na fronteira considerada e ao tempo médio de permanência dessas mesmas unidades. Aplicada à fila de inspeção, a permanência termina quando começa o atendimento. Aplicada à estação inteira, termina na saída da inspeção. Essa distinção permite interpretar uma fila extensa tanto à luz da intensidade das chegadas quanto da duração da espera. Em um modelo de vertiporto com capacidade finita de estacionamento, \citet{onat2026buffer} explicitam outra dependência: quando as posições estão ocupadas, o TLOF fica bloqueado para novas chegadas, que aguardam na fila até a liberação de uma posição.

No lado ar, a quantidade de posições disponíveis limita quantas aeronaves podem receber atendimento simultâneo. A vazão observada é o número de movimentos realizados por hora. Estimar uma capacidade prática exige, além da infraestrutura, declarar o critério de atraso sob o qual essa vazão pode ser sustentada \citep{zhang2023capacity,preisHornung2022heuristic}. Em uma rede estacionária de filas, \citet{escribanoMacias2023} dimensionam pouso, recarga, estacionamento e decolagem pela espera das aeronaves em cada tipo de pad. Sob FCFS, \citet{guerreiro2020} identificam configurações limitadas pelo pouso e outras pelo estacionamento. A reserva técnica após a solicitação difere da admissão antecipada da D2. Com recursos fixos, \citet{halderGandhi2025} já comparam FCFS e reordenação de partidas; medem voos concluídos e ociosidade da frota, sem fila por estação do terminal. A política de sequência também difere da admissão antecipada aqui examinada.

O stand conecta a fila de aeronaves ao processamento dos passageiros. \citet{preisHornung2022} já ligam operações de pad e gate ao atraso agregado do passageiro, inclusive espera no terminal por aeronave atribuída; não discriminam, porém, a fila por estação de inspeção. O conceito operacional descrito por \citet{kunchulia2021} vincula a abertura do gate à prontidão da aeronave. No sentido inverso, o embarque depende da disponibilidade dos passageiros que integrarão o voo \citep{saxena2025vertiminal}. Converter movimentos em passageiros por hora exige ocupação dos assentos e sequência de atendimento; \citet{preis2023} calcula fluxo sob hipóteses de lotação e simultaneidade. Na D2 sintética, a recarga declarada termina antes do acesso à aeronave. A liberação do stand depende da recarga e do último passageiro confirmado que conclui as etapas obrigatórias. Uma posição adicional de estacionamento permite receber outra aeronave enquanto as demais recarregam ou aguardam passageiros, mas não altera as filas do terminal. Já um servidor adicional no balcão atua sobre a prontidão do grupo. \citet{feldhoff2022} observou, sem acoplar o lado ar, que dois balcões reduzem a espera no balcão, e atribuiu a espera baixa na inspeção ao atraso gerado na estação anterior. Como hipótese a ser testada, espera-se que o servidor adicional antecipe o embarque quando o grupo retém a partida. Se a aeronave ainda estiver recarregando, a espera pode apenas se deslocar para a sala de embarque ou para a estação seguinte.

A avaliação do serviço requer observar onde e para quem essa espera ocorre. \citet{feldhoff2022} distingue o tempo total no processamento da parcela de passageiros cuja espera excede o limite de cada estação. Para o LoS temporal, a D2 adota os limites da Tabela 5, reproduzidos da IATA, e não os limites adaptados pela autora para táxis aéreos na Tabela 6. Já \citet{bennaceur2022} consideram a espera esperada entre a chegada prevista ao vertiporto e a partida programada em um problema de alocação de passageiros e voos. Este estudo combina esse indicador de percurso com a espera por estação definida por Feldhoff, mantendo explícitos seus diferentes marcos. Neste estudo, o nível de serviço temporal corresponde à proporção das visitas com espera observada, da entrada na fila ao início do atendimento, que cumprem o limite da estação. A espera entre o fim da inspeção e o início do embarque registra a demora posterior ao processamento, à espera da aeronave ou do grupo. A satisfação percebida exigiria uma medida declarada pelo passageiro.

O resumo de \citet{qin2026commuting} relaciona capacidade do vertiporto, congestionamento e escolha do horário de partida. Em simulação de rede, \citet{onat2024} relacionam tamanho de frota, sob uma política fixa de recarga, ao atraso agregado dos passageiros. Aqui, o recorte é o percurso entre as estações do terminal e a prontidão da aeronave no stand. Esses mecanismos orientam duas comparações. Para investigar atendimento simultâneo, varia-se a quantidade de um recurso com a política e a demanda fixas. Para investigar alocação, compara-se a admissão antecipada por janela e recursos com a alocação local por ordem de solicitação. Infraestrutura, pedidos e requisitos de prontidão são comuns às duas políticas. Em ambos os casos, as decolagens por hora na janela de observação devem ser lidas junto à espera dos passageiros, às recusas e aos atendimentos pendentes. Quando a admissão muda, a composição do grupo que chega a embarcar também pode mudar.
